\makeabstract
    {
    Mercury Contamination due to Legacy Pollution in Puffer{\textquoteright}s Pond, Amherst, MA
    }
    {
    Anika Ponnambalam\textsuperscript{1}, Alysha Rodr{\'\i}guez Cab\'an\textsuperscript{1}, Anna Martini\textsuperscript{1}
    }
    {
    \textsuperscript{1} Department of Geology, Amherst College, Amherst, MA
    }
    {
    The Town of Amherst is planning to dredge Puffer{\textquoteright}s Pond as part of a larger improvement project. Occasionally, Puffer{\textquoteright}s Pond is closed due to high E. coli counts. Dredging can improve water quality, which in turn can lower the amount of E. coli; however, there is a known interval of mercury-laden sediments in Puffer{\textquoteright}s Pond resulting from legacy pollution. Thus, this project seeks to study the potential negative environmental impacts of resuspending the mercury-laden sediment to assess future action.
    }
    {
    Local fish were caught via lure from Puffer{\textquoteright}s Pond and Hadley Reservoir, a spring fed waterbody for comparison. Fish were measured and dissected, separating liver and muscle samples. Fish samples were then weighed to obtain wet mass, freeze dried, and weighed again to obtain dry mass. The dried samples were run on a Nippon direct mercury analyzer. To obtain wet mercury values, the percent water content in each specimen was calculated using the wet and dry mass. Then, that percentage was used to calculate the wet Hg using the formula: Wet Hg = 100 - (\%H2O) * 0.01 * (Dry Hg)
    }
    {
    Non-stocked fish samples from Puffer{\textquoteright}s Pond had average wet Hg concentrations of 1180.20, 47.89, and 203.88 ppb, while stocked samples had concentrations of 27.06 and 21.64 ppb. In comparison to Hadley Reservoir, there was no statistically significant difference between fish from the two sites due to the small sample size. However, the specimen with 1180.20 ppb from Puffer{\textquoteright}s Pond had a distribution of Hg in its tissues indicative of significant contamination: its liver:muscle Hg ratio was \ensuremath{>}1, meaning that its liver was likely working to detoxify the body due to high Hg levels (Cizdziel et al. 2013).
    }
    {
    Dredging Puffer{\textquoteright}s Pond could resuspend the layer of mercury-rich sediment and increase Hg contamination in the water. This creates a question for future research: If Puffer{\textquoteright}s Pond is dredged, will wet mercury concentration in fish increase relative to fish in a nearby uncontaminated environment (i.e. Hadley Reservoir)? Further research is needed to establish a more significant baseline for future comparison.
    }

    \newpage
    